% ================================================================
% ==== FILE: content/cycles/cycles_k7.tex
% ================================================================

\subsubsection{Zyklen auf \texorpdfstring{$K_7$}{K_7}}
\label{sec:cycles-k7}

Das Niveau$K_7$stellt das soziale Kontinuum dar: eine strukturierte Domäne von
Rollen, Normen, Kommunikationsflüsse, Vertrauensdynamik und institutionelle Aspekte
Stabilisierung.
Zyklen weiter$K_7$Formalisierung wiederkehrender sozialer Prozesse, die das aufrechterhalten
Beharrlichkeit sozialer Gruppen, Institutionen und kollektiven Verhaltens.

Äxte an$K_7$enthalten:
\[
  A_{\mathrm{role}},\;
  A_{\mathrm{norm}},\;
  A_{\mathrm{trust}},\;
  A_{\mathrm{comm}},\;
  A_{\mathrm{coop}},
\]
mit Potenzialen:
\[
  P_{\mathrm{trust}},\;
  P_{\mathrm{coh}},\;
  P_{\mathrm{normative}},\;
  P_{\mathrm{stability}},
\]
and flows (the corresponding derivation note):
\[
  J_{\mathrm{comm}},\;
  J_{\mathrm{coop}},\;
  J_{\mathrm{stab}}.
\]

Schwellenwerte, die die Lebensfähigkeit des Zyklus bestimmen:
\[
  \Theta_{\mathrm{trust}},\;
  \Theta_{\mathrm{norm}},\;
  \Theta_{\mathrm{coh\text{-}soc}},\;
  \Theta_{\mathrm{embed}},\;
  \Theta_{\mathrm{fragility}}.
\]

\subsubsection{Überblick}

Ein soziales Kontinuum besteht nicht, weil seine Komponenten statisch sind, sondern weil
sie werden durch wiederkehrende Kommunikationszyklen dynamisch aufrechterhalten,
Zusammenarbeit, Durchsetzung von Normen, Wiederherstellung des Vertrauens und institutionelle Aspekte
Reproduktion.
In diesem Abschnitt werden diese Zyklen in der Sprache des OC-Frameworks formalisiert.

\subsubsection{1. Kommunikationszyklus$C_{\mathrm{comm}}$}

Kommunikation ist der Kernprozess zur Schaffung eines gemeinsamen Zustands innerhalb des Sozialen
Kontinuum.

Der Zyklus:
\[
  \mathrm{signal}
  \to
  \mathrm{interpretation}
  \to
  \mathrm{feedback}
  \to
  \mathrm{alignment}
  \to
  \mathrm{signal}.
\]

Formal:
\[
  C_{\mathrm{comm}}
    =
    \left\{
      A_{\mathrm{comm}}(t),
      P_{\mathrm{coh}}(t),
      J_{\mathrm{comm}}(t)
    \right\}.
\]

Stabilität erfordert:
\[
  T_{\mathrm{comm}} < \Theta_{\mathrm{coh\text{-}soc}}.
\]

\subsubsection{2. Vertrauens-Regenerationszyklus$C_{\mathrm{trust}}$}

Trust is a dynamic potential $P_{\mathrm{trust}}$ subject to fluctuations,
Erosion und Wiederherstellung.

Zyklusstruktur:
\[
  \mathrm{cooperation}
  \to
  \mathrm{verification}
  \to
  \mathrm{reinforcement}
  \to
  \mathrm{renewal}
  \to
  \mathrm{cooperation}.
\]

Formal:
\[
 C_{\mathrm{trust}}
   =
   \left\{
     P_{\mathrm{trust}}(t),\;
     J_{\mathrm{coop}}(t),\;
     A_{\mathrm{trust}}(t)
   \right\}.
\]

Existenzbedingung:
\[
  P_{\mathrm{trust}}(t+\tau_{\mathrm{trust}})=
  P_{\mathrm{trust}}(t).
\]

Schwelle:
\[
  T_{\mathrm{trust}} < \Theta_{\mathrm{trust}}.
\]

\subsubsection{3. Normdurchsetzungszyklus$C_{\mathrm{norm}}$}

Soziale Normen regeln Erwartungen und Verhalten.

Zyklus:
\[
  \mathrm{activation}
  \to
  \mathrm{monitoring}
  \to
  \mathrm{sanction}
  \to
  \mathrm{repair}
  \to
  \mathrm{activation}.
\]

Bilden:
\[
  C_{\mathrm{norm}}=
  \left\{
    A_{\mathrm{norm}}(t),
    P_{\mathrm{normative}}(t),
    J_{\mathrm{stab}}(t)
  \right\}.
\]

Schwellenwertbedingung:
\[
  T_{\mathrm{norm}} < \Theta_{\mathrm{norm}}.
\]

\subsubsection{4. Rollenstrukturzyklus$C_{\mathrm{role}}$}

Soziale Rollen müssen kontinuierlich reproduziert werden, um die institutionelle Ordnung aufrechtzuerhalten.

Zyklus:
\[
  \mathrm{role\ learning}
  \to
  \mathrm{performance}
  \to
  \mathrm{evaluation}
  \to
  \mathrm{adjustment}
  \to
  \mathrm{role\ learning}.
\]

Formal:
\[
  C_{\mathrm{role}}
  =
  \left\{
     A_{\mathrm{role}}(t),\;
     P_{\mathrm{stability}}(t),\;
     J_{\mathrm{comm}}(t)
  \right\}.
\]

Schwelle:
\[
  T_{\mathrm{role}} < \Theta_{\mathrm{coh\text{-}soc}}.
\]

\subsubsection{5. Kooperationszyklus$C_{\mathrm{coop}}$}

Cooperation, sustained by $J_{\mathrm{coop}}$, is a recurrent structure:

\[
 \mathrm{proposal}
 \to
 \mathrm{coordination}
 \to
 \mathrm{execution}
 \to
 \mathrm{reward}
 \to
 \mathrm{renewal}.
\]

Formal:
\[
 C_{\mathrm{coop}}=
 \left\{
   J_{\mathrm{coop}}(t),\;
   P_{\mathrm{trust}}(t),\;
   A_{\mathrm{coop}}(t)
 \right\}.
\]

Stabilitätsbedingung:
\[
 T_{\mathrm{coop}} < \Theta_{\mathrm{trust}}.
\]

\subsubsection{6. Institutioneller Reproduktionszyklus$C_{\mathrm{inst}}$}

Institutionen sind dauerhafte soziale Strukturen, die definiert sind durch:

\[
  \mathrm{codification}
  \to
  \mathrm{implementation}
  \to
  \mathrm{monitoring}
  \to
  \mathrm{repair}
  \to
  \mathrm{codification}.
\]

Die institutionelle Reproduktion hängt ab von:

\[
 C_{\mathrm{inst}}
 =
 \left\{
   A_{\mathrm{norm}}(t),\;
   A_{\mathrm{role}}(t),\;
   P_{\mathrm{stability}}(t),\;
   J_{\mathrm{stab}}(t)
 \right\}.
\]

Schwelle:
\[
 T_{\mathrm{inst}} < \Theta_{\mathrm{embed}}.
\]

Die Einbettungsraumbeschränkung kodiert die Abhängigkeit von
$K_7$An$M_7$.

\subsubsection{7. Sozialer Stabilitätszyklus$C_{\mathrm{stab}}$}

Dies ist eine mehrzyklische Integration:

\[
  C_{\mathrm{stab}}
    =
    C_{\mathrm{comm}}
    \cup
    C_{\mathrm{trust}}
    \cup
    C_{\mathrm{norm}}
    \cup
    C_{\mathrm{role}}
    \cup
    C_{\mathrm{coop}}.
\]

Stabilität erfordert:

\[
  T_{\mathrm{soc}}
  <
  \min(\Theta_{\mathrm{trust}},
       \Theta_{\mathrm{norm}},
       \Theta_{\mathrm{coh\text{-}soc}},
       \Theta_{\mathrm{embed}}).
\]

Dieser Zyklus entspricht der Existenz eines kohärenten sozialen Kontinuums.

\subsubsection{Zyklusmetriken für \texorpdfstring{$K_7$}{K_7}}

\paragraph{Length.}
\[
L(C)
  = \oint d\Omega(K_7).
\]

\paragraph{Efficiency.}
\[
  C_{\mathrm{eff}}
    =
    \frac{\oint J_7\cdot dA}{L(C)}.
\]

\paragraph{Stability.}
\[
SC)
  = \min_{t\in C}
    d(\Omega(K_7), \partial\Omega(K_7)).
\]

\paragraph{Weight.}
\[
WC)
 =
 F(C_{\mathrm{eff}},S(C),
   P_{\mathrm{trust}},
   P_{\mathrm{normative}},
   P_{\mathrm{coh}}).
\]

\subsubsection{Zusammenbruch von \texorpdfstring{$K_7$}{K_7} Zyklen}

Ein Zusammenbruch sozialer Zyklen tritt auf, wenn:

\[
 T_{\mathrm{trust}} > \Theta_{\mathrm{trust}},
 \qquad
 T_{\mathrm{norm}} > \Theta_{\mathrm{norm}},
 \qquad
 T_{\mathrm{comm}} > \Theta_{\mathrm{coh\text{-}soc}},
\]

oder wenn Einschränkungen hinsichtlich des Einbettungsraums fehlschlagen:

\[
 M_7 \not\supseteq \Omega(K_7).
\]

Zusammenbruch bedeutet:

\[
 C_j \to \varnothing,\quad
 \Omega(K_7)=\varnothing,\quad
 k_7\to 0.
\]

\subsubsection{Kontinuität von \texorpdfstring{$K_6$}{K_6} zu \texorpdfstring{$K_7$}{K_7}}

The transition $K_6\to K_7$ (the corresponding derivation note) occurs when:
\[
  T_{\mathrm{comm}} > \Theta_{\mathrm{dim}},
\]
und neue Achsen erscheinen:
\[
  A_{\mathrm{role}},\;
  A_{\mathrm{norm}},\;
  A_{\mathrm{trust}},\;
  A_{\mathrm{coop}}.
\]

Dies schafft neue Schwellen und ermöglicht die Existenz sozialer Kreisläufe
ohne Analogie auf kognitiver Ebene$K_6$.
